\section{Research Design}

This study, ``FeatureSAGE: Stable Attribute Group Editing With Integrated StyleFeatureEditing Encoder for Enhanced Few-Shot Image Generation,'' 
adopts a quantitative experimental research design to evaluate improvements to the existing SAGE framework. In particular, the GAN inversion 
module of SAGE is modified by integrating the StyleFeatureEditor (SFE) inverter and its Feature Editor module. The modification is evaluated 
under a one-shot setting, where attribute manipulation is conditioned on a single reference image per test instance. Objective evaluation is 
performed using Fr\'echet Inception Distance (FID) and Learned Perceptual Image Patch Similarity (LPIPS), which measure distributional image 
quality and perceptual variation, respectively.

% =============================================================================
\subsection{Experimental Setup}

The experimental setup is designed to evaluate the effectiveness of the proposed FeatureSAGE framework in comparison 
to the existing SAGE method. The evaluation follows a quantitative experimental design that involves controlled 
changes to the inversion and feature-editing components and the measurement of model outputs using standardized metrics.

\begin{itemize}
    \item \textbf{Independent Variable:} The main experimental manipulation is the GAN inversion method. The baseline SAGE configuration uses the original pixel2style2pixel (pSp) encoder, whereas the proposed FeatureSAGE configuration integrates the StyleFeatureEditor inverter and Feature Editor module.

    \item \textbf{Dependent Variables:}
    To objectively evaluate the effectiveness of the GAN inversion and editing pipeline, the following metrics are employed:
    \begin{itemize}
        \item Reconstruction quality and generated-image distribution alignment, measured using Fr\'echet Inception Distance (FID).
        \item Perceptual similarity or perceptual diversity, depending on the evaluation context, measured using Learned Perceptual Image Patch Similarity (LPIPS).
    \end{itemize}

    \item \textbf{Controlled Variables}
    \begin{enumerate}
        \item \textbf{Generator Architecture:} A fixed, pretrained StyleGAN2 generator is used in both configurations for consistent image synthesis capability.
        \item \textbf{Training Dataset:} Both configurations use the same FUNIT Animal Faces and 102 Flowers datasets with the same preprocessing and seen/unseen class splits.
        \item \textbf{Evaluation Protocol:} The same unseen categories, one-shot support samples, query sets, and evaluation procedures are used for both configurations.
        \item \textbf{Loss Functions:} Reconstruction and editing losses are applied consistently within the corresponding training phases.
        \item \textbf{Optimization Settings:} Both configurations use the Adam optimizer and the fixed learning-rate settings specified in the training strategy.
        \item \textbf{Evaluation Metrics:} FID and LPIPS are computed using the same implementations and parameters.
        \item \textbf{Hardware and Software Environment:} All experiments are conducted using the same GPU-enabled PyTorch environment to avoid performance inconsistencies.
    \end{enumerate}

    \item \textbf{Experimental Conditions}
    
    Two experimental configurations are used:
    \begin{itemize}
        \item SAGE (Baseline): Employs the original pSp encoder with StyleGAN2.
        \item FeatureSAGE (Proposed): Integrates the StyleFeatureEditor inverter and Feature Editor.
    \end{itemize}
\end{itemize}

% =========================================================================================

\subsection{Hardware and System Requirements}

All experiments were conducted on a virtual machine instance with GPU pass-through hosted on institutional cloud infrastructure. The same hardware and software configuration was used for the baseline and proposed models so that performance differences could be attributed to the inversion and feature-editing components rather than to changes in the execution environment. Table~\ref{tab:hardware_system_requirements} summarizes the system requirements reported for the implementation.

\begin{table}[H]
\centering
\setlength{\tabcolsep}{8pt}
\begin{tabular}{|p{0.28\linewidth}|p{0.62\linewidth}|}
\hline
\textbf{Component} & \textbf{Specification} \\ \hline
CPU & QEMU Virtual CPU version 2.5+, 2 processors at 2.19 GHz \\ \hline
RAM & 32.0 GB \\ \hline
GPU & NVIDIA GRID V100D-32A, virtualized Tesla V100 with 16 GB HBM2 \\ \hline
Storage & Not specified in the source implementation notes; storage is required for the datasets, model checkpoints, and generated outputs. \\ \hline
Software and frameworks & PyTorch environment; pretrained StyleGAN2 generator; pSp encoder for the baseline and editing-pair generation; StyleFeatureEditor inverter and Feature Editor; MTCNN for animal-face cropping and alignment; Adam optimizer; mixed precision for Feature Editor training; FID and LPIPS metric implementations. \\ \hline
\end{tabular}
\caption{Hardware and system requirements used in the experimental setup.}
\label{tab:hardware_system_requirements}
\end{table}

The GPU specification is important because FeatureSAGE performs multiple generator passes and feature-tensor operations during inversion and editing. The use of mixed precision in the Feature Editor phase reduces memory pressure on the 16 GB V100 vGPU while preserving the same model configuration and reported evaluation protocol.

% =========================================================================================

\subsection{Few-Shot Experimental Setting}

In this study, the term \emph{few-shot} refers to the number of reference images available per target class used to construct 
class-level latent representations for attribute group editing and downstream evaluation. A \textbf{one-shot} setting utilizes a single reference image 
per target class, representing an extreme data-scarce scenario in which attribute editing and inversion must be performed with minimal visual 
information. This one-shot configuration is applied consistently across all datasets and experimental conditions for both the baseline SAGE 
and the proposed FeatureSAGE framework. By explicitly controlling the evaluation to a single reference image per class, the experimental 
protocol enables a systematic assessment of FeatureSAGE's robustness, stability, and generalization capability under severe data limitations. 
The few-shot constraint is applied during test-time adaptation and evaluation on unseen classes, while the StyleGAN2 generator and attribute 
directions are learned from seen-category data.
